% GENERATED FILE. Edit canonical lessons, then run npm run generate:books.
% canonical-source-hash: fnv1a64:63c4f816b923b079
% canonical-lessons: ES-C34-pensar, ES-C34-entender, ES-C34-leer, ES-C34-escribir

\chapter{Four Verbs of the Mind}
\label{ch:spanish-verbs-of-the-mind}
\hlchaptermodality{41}

\begin{chapteropening}
\textbf{By the end of this chapter:} \emph{I can say what I think, understand, read and write in Spanish, and break the stems of pensar and entender the way a stem-changer requires.}

Run pensar and entender through the e$\to$ie boot, then leer and escribir as the pair that does not break, and close on the propped-up e- that turns scrībere into escribir and schola into escuela.
\end{chapteropening}

\section[pensar]{pensar — \textquotedblleft{}to think\textquotedblright{}, and it started out meaning \textquotedblleft{}to weigh\textquotedblright{}}

\label{lesson:ES-C34-pensar}



\begin{quote}
You can already run an \emph{-ar} verb: \emph{hablo, hablas, habla}. Here is one that takes those same endings and then cracks its stem in the middle, the way \emph{querer} did.
\end{quote}

\begin{sounds}
\begin{itemize}
\raggedright
  \item \texttt{diphthong-ie} — under stress the \emph{e} breaks open: \textbf{pienso} = \emph{PYEN-so}.
  \item The \emph{s} is always a clear hissing \emph{s}, never an English \emph{z} buzz.
\end{itemize}
\end{sounds}

\begin{cousinweb}
\textbf{pensar} (\textquotedblleft{}\textbf{to think}\textquotedblright{}) $\leftarrow$ Latin \textbf{pēnsāre}, \textquotedblleft{}\textbf{to weigh}\textquotedblright{} — the repeating form of \emph{pendere}, \textquotedblleft{}to hang, to weigh out\textquotedblright{}. A Roman hung goods on a balance; the same verb slid across to weighing an idea.

English kept the whole family, and most of it still smells of money and scales:

\begin{itemize}
\raggedright
  \item \textbf{pensive} — turning something over on the balance.
  \item \textbf{compensate} — \emph{com-} \textquotedblleft{}together\textquotedblright{} + weigh: settle one weight against another.
  \item \textbf{expense}, \textbf{dispense}, \textbf{pension} — payments were literally \emph{weighed out}.
  \item \textbf{peso} — a coin whose name simply means \textquotedblleft{}a weight\textquotedblright{}.
  \item \textbf{pansy} — from French \emph{pensée}, \textquotedblleft{}a thought\textquotedblright{}.
\end{itemize}

English reaches for the same picture without noticing it: you \emph{weigh your options}, and an argument \emph{carries weight}.
\end{cousinweb}

\begin{grammarlens}[title={Grammar Lens: the same endings, a broken stem}]
The endings are the ordinary \emph{-ar} set. Only the stressed \emph{e} changes:

\noindent
\begin{tabularx}{\linewidth}{@{}>{\raggedright\arraybackslash}X>{\raggedright\arraybackslash}X>{\raggedright\arraybackslash}X@{}}
\toprule
\textbf{person} & \textbf{form} & \textbf{stem} \\
\midrule
yo & \textbf{pienso} & broken \\
tú & \textbf{piensas} & broken \\
él/ella/usted & \textbf{piensa} & broken \\
nosotros & \textbf{pensamos} & plain \\
ellos/ustedes & \textbf{piensan} & broken \\
\bottomrule
\end{tabularx}

\emph{Pensamos} keeps the plain \emph{e} because there the stress lands on the ending, not the stem — exactly the rule \emph{querer} taught you.

Two small habits come with the verb: \emph{pensar} \textbf{en} something means \textquotedblleft{}think about\textquotedblright{} it (\emph{pienso en ti}), and \emph{pensar} followed by a bare infinitive means \textquotedblleft{}plan to\textquotedblright{} (\emph{pienso hablar} — I mean to speak).
\end{grammarlens}

\subsection*{Your turn}
Say these aloud:

\begin{itemize}
\raggedright
  \item \textquotedblleft{}pensar\textquotedblright{} then the set — \textquotedblleft{}pienso, piensas, piensa, pensamos, piensan\textquotedblright{}
  \item \textquotedblleft{}Pienso en ti\textquotedblright{} — I am thinking of you
  \item \textquotedblleft{}pensar\textquotedblright{} then English \textquotedblleft{}pensive, compensate, peso\textquotedblright{} — the weighing family
\end{itemize}

\begin{tcolorbox}[breakable,colback=teal!4,colframe=teal!35!black,title={Before you move on}]
What did \emph{pēnsāre} mean literally, and name two English cousins. (\textbf{To weigh}; \emph{pensive, compensate, expense, peso}.) Which form keeps the plain \emph{e}, and why? (\emph{Pensamos} — the stress falls on the ending.) Which little word follows \emph{pensar} for \textquotedblleft{}think about\textquotedblright{}? (\textbf{En}.) Next: \textbf{entender}, which cracks the same way.
\end{tcolorbox}
\section[entender]{entender — \textquotedblleft{}to understand\textquotedblright{}, built out of stretching}

\label{lesson:ES-C34-entender}



\begin{quote}
\emph{Pensar} broke its stem under stress. So does this one — and it hands you the single most useful sentence a beginner owns: \textbf{no entiendo}, \textquotedblleft{}I don't understand.\textquotedblright{}
\end{quote}

\begin{sounds}
\begin{itemize}
\raggedright
  \item \texttt{diphthong-ie} — \textbf{entiendo} = \emph{en-TYEN-do}.
  \item \texttt{d-soft} — the \emph{d} between vowels is soft, close to the \emph{th} of English \emph{this}.
\end{itemize}
\end{sounds}

\begin{cousinweb}
\textbf{entender} (\textquotedblleft{}\textbf{to understand}\textquotedblright{}) $\leftarrow$ Latin \textbf{intendere}: \textbf{in-} (\textquotedblleft{}toward\textquotedblright{}) + \textbf{tendere} (\textquotedblleft{}\textbf{to stretch}\textquotedblright{}). Understanding something was \emph{stretching your attention toward} it — leaning in.

\emph{tendere} stretches right across English:

\begin{itemize}
\raggedright
  \item \textbf{intend}, \textbf{intent}, \textbf{intense}, \textbf{intention} — the very same \emph{in-} + stretch.
  \item \textbf{attend}, \textbf{attention} — \emph{ad-} \textquotedblleft{}toward\textquotedblright{}, stretch that way.
  \item \textbf{extend}, \textbf{contend}, \textbf{pretend}, \textbf{tension}, \textbf{tendon}, and \textbf{tent}, which is stretched cloth.
\end{itemize}

So \emph{entiendo} and \emph{I intend} are the same two Latin pieces landing in different places: Spanish went to \textquotedblleft{}I grasp it\textquotedblright{}, English to \textquotedblleft{}I mean to\textquotedblright{}.

\begin{quote}
\textbf{Careful with a false friend.} Spanish \textbf{pretender} is \emph{not} English \emph{pretend}. It means \textquotedblleft{}to aim at, to be after\textquotedblright{} — \emph{pretende ser médico}, \textquotedblleft{}he means to become a doctor\textquotedblright{}. Same pieces, different landing.
\end{quote}
\end{cousinweb}

\begin{grammarlens}[title={Grammar Lens: an -er verb that cracks, and a calmer twin}]
Ordinary \emph{-er} endings, stressed \emph{e} broken to \emph{ie}:

\noindent
\begin{tabularx}{\linewidth}{@{}>{\raggedright\arraybackslash}X>{\raggedright\arraybackslash}X>{\raggedright\arraybackslash}X@{}}
\toprule
\textbf{person} & \textbf{form} & \textbf{stem} \\
\midrule
yo & \textbf{entiendo} & broken \\
tú & \textbf{entiendes} & broken \\
nosotros & \textbf{entendemos} & plain \\
ellos/ustedes & \textbf{entienden} & broken \\
\bottomrule
\end{tabularx}

Spanish has a second verb for this job: \textbf{comprender} $\leftarrow$ Latin \emph{comprehendere}, \textquotedblleft{}to seize together\textquotedblright{} ($\to$ English \textbf{comprehend}, \textbf{apprehend}, \textbf{comprehensive}). It is a plain \emph{-er} verb with no crack at all: \emph{comprendo, comprendes, comprende}. Both mean \textquotedblleft{}understand\textquotedblright{}. \emph{Entender} is the everyday one; \emph{comprender} leans a little more formal, and more toward grasping a whole thing than catching a word.
\end{grammarlens}

\subsection*{Your turn}
Say these aloud:

\begin{itemize}
\raggedright
  \item \textquotedblleft{}No entiendo\textquotedblright{} — I don't understand
  \item \textquotedblleft{}entiendo, entiendes, entendemos, entienden\textquotedblright{}
  \item \textquotedblleft{}pienso\textquotedblright{} then \textquotedblleft{}entiendo\textquotedblright{} — the same crack, twice
  \item \textquotedblleft{}entender\textquotedblright{} then English \textquotedblleft{}intend, intense, attention\textquotedblright{} — the stretching family
\end{itemize}

\begin{tcolorbox}[breakable,colback=teal!4,colframe=teal!35!black,title={Before you move on}]
What does \emph{tendere} mean, and what does \emph{in-} add? (\textbf{To stretch}; \textquotedblleft{}toward\textquotedblright{}.) Give the \emph{yo} and \emph{nosotros} forms. (\emph{Entiendo, entendemos}.) Of \emph{entender} and \emph{comprender}, which one breaks its stem? (\textbf{Entender}.) And what does Spanish \emph{pretender} mean? (\textbf{To aim at}, not to pretend.) Next: \textbf{leer}.
\end{tcolorbox}
\section[leer]{leer — \textquotedblleft{}to read\textquotedblright{}, and before that, \textquotedblleft{}to gather\textquotedblright{}}

\label{lesson:ES-C34-leer}



\begin{quote}
Two cracked verbs down. Here is a calm one — nothing breaks in the middle. The only difficulty is in your mouth: two \emph{e}'s in a row that must stay two.
\end{quote}

\begin{sounds}
\begin{itemize}
\raggedright
  \item \texttt{hiatus-ee} — \textbf{leer} is two syllables, \emph{le-ER}, never one long English \emph{ee}.
  \item \textbf{leo} is likewise two beats, \emph{LE-o}.
\end{itemize}
\end{sounds}

\begin{cousinweb}
\textbf{leer} (\textquotedblleft{}\textbf{to read}\textquotedblright{}) $\leftarrow$ Latin \textbf{legere}, \textquotedblleft{}\textbf{to gather, to pick out, to choose}\textquotedblright{} — and only afterwards \textquotedblleft{}to read\textquotedblright{}. A reader picks letters out of a row the way a farmhand picks olives off a branch.

That gathering root runs through an enormous amount of English:

\begin{itemize}
\raggedright
  \item \textbf{legible}, \textbf{legend} (\textquotedblleft{}things to be read\textquotedblright{}), \textbf{lecture}, \textbf{lectern}, and worn right down, \textbf{lesson}.
  \item built with prefixes: \textbf{collect} (\emph{col-} \textquotedblleft{}together\textquotedblright{}), \textbf{select} (\emph{se-} \textquotedblleft{}apart\textquotedblright{}), \textbf{elect} (\emph{e-} \textquotedblleft{}out\textquotedblright{}), \textbf{intellect} (\emph{inter-} \textquotedblleft{}between\textquotedblright{} — choosing between things).
  \item \textbf{diligent} and \textbf{elegant}, both of which meant \textquotedblleft{}picking carefully\textquotedblright{}.
\end{itemize}

One honest limit: Greek \emph{légein} (\textquotedblleft{}to speak, to gather\textquotedblright{}) — which gives English \emph{logic}, \emph{dialogue}, \emph{catalogue} — is a \textbf{cousin} of \emph{legere} from the same ancient root, not its parent. Same family, separate journeys.
\end{cousinweb}

\begin{grammarlens}[title={Grammar Lens: a plain -er verb, with a mouthful}]
Nothing irregular happens. Drop \emph{-er}, add the ordinary endings:

\noindent
\begin{tabularx}{\linewidth}{@{}>{\raggedright\arraybackslash}X>{\raggedright\arraybackslash}X>{\raggedright\arraybackslash}X@{}}
\toprule
\textbf{person} & \textbf{form} & \textbf{said} \\
\midrule
yo & \textbf{leo} & LE-o \\
tú & \textbf{lees} & LE-es \\
él/ella/usted & \textbf{lee} & LE-e \\
nosotros & \textbf{leemos} & le-E-mos \\
ellos/ustedes & \textbf{leen} & LE-en \\
\bottomrule
\end{tabularx}

Every one of those is a syllable longer than an English speaker expects. Give each \emph{e} its own beat and the word comes out right.
\end{grammarlens}

\subsection*{Your turn}
Say these aloud:

\begin{itemize}
\raggedright
  \item \textquotedblleft{}leer\textquotedblright{} — \emph{le-ER}, two beats
  \item \textquotedblleft{}Leo un libro\textquotedblright{} — I read a book
  \item \textquotedblleft{}No entiendo\textquotedblright{} then \textquotedblleft{}Leo, pero no entiendo\textquotedblright{} — I read, but I don't understand
  \item \textquotedblleft{}leer\textquotedblright{} then English \textquotedblleft{}legible, collect, select, elect\textquotedblright{}
\end{itemize}

\begin{tcolorbox}[breakable,colback=teal!4,colframe=teal!35!black,title={Before you move on}]
What did \emph{legere} mean before it meant \textquotedblleft{}to read\textquotedblright{}? (\textbf{To gather, to pick out}.) Give two English words built on it with a prefix. (\emph{Collect, select, elect, intellect}.) How many syllables in \emph{leo}? (\textbf{Two} — \emph{LE-o}.) Next: \textbf{escribir}, the other half of literacy.
\end{tcolorbox}
\section[escribir]{escribir — \textquotedblleft{}to write\textquotedblright{}, which once meant \textquotedblleft{}to scratch\textquotedblright{}}

\label{lesson:ES-C34-escribir}



\begin{quote}
\emph{Leer} gathers letters off a page. This is the other half of the same skill — and its root remembers a time when writing was done with a sharp point in wax.
\end{quote}

\begin{sounds}
\begin{itemize}
\raggedright
  \item \texttt{st-no-e} — Spanish will not begin a word with \emph{sc-}, \emph{sp-}, \emph{st-}, so it props one up with an \textbf{e-}: \emph{es-cri-BIR}.
  \item \texttt{b-soft} — between vowels the \emph{b} is soft, lips barely touching.
\end{itemize}
\end{sounds}

\begin{cousinweb}
\textbf{escribir} (\textquotedblleft{}\textbf{to write}\textquotedblright{}) $\leftarrow$ Latin \textbf{scrībere}, whose first sense was \textquotedblleft{}\textbf{to scratch, to incise}\textquotedblright{}. Writing was cutting a line into wax or stone before it was ever ink on paper.

The English family is huge, and every member still works with a pen:

\begin{itemize}
\raggedright
  \item \textbf{scribe}, \textbf{scribble}, \textbf{script}, \textbf{scripture}.
  \item \textbf{describe} (\emph{de-} \textquotedblleft{}down\textquotedblright{}) and \textbf{description}; \textbf{prescribe} (\emph{prae-} \textquotedblleft{}before\textquotedblright{}) and \textbf{prescription}; \textbf{subscribe} (\emph{sub-} \textquotedblleft{}under\textquotedblright{} — you signed at the bottom); \textbf{inscription}; \textbf{transcript}; \textbf{postscript}.
  \item \textbf{manuscript} — \emph{manus} \textquotedblleft{}hand\textquotedblright{} + \emph{scriptum}: written by hand.
  \item \textbf{conscription} — being \emph{written into} an army.
\end{itemize}

The Spanish past participle is \textbf{escrito} (\textquotedblleft{}written\textquotedblright{}), which is English \textbf{script} wearing different clothes.
\end{cousinweb}

\begin{grammarlens}[title={Grammar Lens: an -ir verb, and a productive spelling rule}]
\emph{Escribir} is a plain \emph{-ir} verb, the family \emph{vivir} opened:

\noindent
\begin{tabularx}{\linewidth}{@{}>{\raggedright\arraybackslash}X>{\raggedright\arraybackslash}X>{\raggedright\arraybackslash}X@{}}
\toprule
\textbf{person} & \textbf{form} & \textbf{ending} \\
\midrule
yo & \textbf{escribo} & -o \\
tú & \textbf{escribes} & -es \\
él/ella/usted & \textbf{escribe} & -e \\
nosotros & \textbf{escribimos} & -imos \\
ellos/ustedes & \textbf{escriben} & -en \\
\bottomrule
\end{tabularx}

That propped-up \textbf{e-} is a rule you can run yourself, not a quirk of one word: Latin \emph{schola} became \textbf{escuela}, \emph{spatha} became \textbf{espada} (sword), \emph{stāre} became \textbf{estar}, \emph{Hispānia} became \textbf{España}. It is also why a Spanish speaker learning English says \emph{eh-Spain}.

Now you can state a plan out loud: \emph{Pienso escribir} — \textquotedblleft{}I mean to write.\textquotedblright{}
\end{grammarlens}

\subsection*{Your turn}
Say these aloud:

\begin{itemize}
\raggedright
  \item \textquotedblleft{}escribo, escribes, escribe, escribimos, escriben\textquotedblright{}
  \item \textquotedblleft{}Leo y escribo español\textquotedblright{} — I read and write Spanish
  \item \textquotedblleft{}Pienso escribir\textquotedblright{} — I mean to write
  \item \textquotedblleft{}escribir\textquotedblright{} then English \textquotedblleft{}scribe, describe, manuscript\textquotedblright{} — the scratching family
\end{itemize}

\begin{tcolorbox}[breakable,colback=teal!4,colframe=teal!35!black,title={Before you move on}]
What did \emph{scrībere} mean first? (\textbf{To scratch, to incise}.) Why does Spanish say \emph{escuela} and \emph{España}? (It refuses to open a word with \emph{sc-} or \emph{sp-}, so it props an \textbf{e-} in front.) Say \textquotedblleft{}I read and write Spanish\textquotedblright{}. (\emph{Leo y escribo español}.) And what does \emph{legere} mean, from the lesson before? (\textbf{To gather, to pick out}.) Next chapter: taking, asking, helping — and a verb that runs backwards.
\end{tcolorbox}
